\documentclass[12pt]{article}
\usepackage[utf8]{inputenc}
\usepackage[T1]{fontenc}
\usepackage{hyperref}
\title{Wagdi Language Resources for K-12 Curriculum Development}
\author{Compiled by Claude Sonnet 5 (claude-sonnet-5) for Andr\'as Kornai}
\date{}
\begin{document}
\maketitle

\section{Introduction}
Wagdi (Vagdi; ISO 639-3 wbr; Glottocode wagd1238) is spoken mainly in the Dungarpur and Banswara districts of southern Rajasthan (the Vagad region). It is usually treated as a Bhil language. Speaker counts in the search results vary: 2.5 million for the 2001 census and 3.39 million for the 2011 census. The dialects listed are Aspur, Kherwara, Sagwara and Adivasi Wagdi. The search found only encyclopedic and academic descriptions.

\section{Primary Educational Resources}
No Wagdi primers, textbooks or school programmes were found in search results.

\section{Linguistic and Reference Materials}
\begin{itemize}
\item Wagdi, Wikipedia: \url{https://en.wikipedia.org/wiki/Wagdi}
\item Dungarpur district, Wikipedia (regional context): \url{https://en.wikipedia.org/wiki/Dungarpur_district}
\item Banswara district, Wikipedia (regional context): \url{https://en.wikipedia.org/wiki/Banswara_district}
\end{itemize}
A 2012 SOAS PhD thesis, ``Dialect Continuum in the Bhil Tribal Belt: Grammatical Aspects'' (doi 10.25501/SOAS.00014048), was mentioned in the search summary. I did not see a direct link to it, so it is not given a URL here.

\section{Religious and Community Resources}
None found.

\section{Digital and Audio Resources}
None found.

\section{Curriculum Development Context}
The language is not known to be taught in schools. Materials would have to be built from scratch, drawing on the SOAS thesis for dialect and grammar information. Hindi is the regional school language.
\end{document}